\scanpage{23}
% archive foliation: 22
\uline{En résumé} Dans la \uline{catégorie}
\uline{abélienne} \inserted{$k$-lin. $C$} de tous les $k$-modules
$\frac12$ simpliciaux, ceux qui sont \inserted{str. $\partial$-} parfaits
forment une sous-catégorie \uline{additive}
$C_0$ \unsure{en question|engendrée}, dans laquelle
on a un complexe de \unsure{cochaînes}
\diagram{\begin{tikzcd}
0\arrow[r,dash]&\Phi_*^{(0)}\arrow[r,dash]&\Phi_*^{(1)}\arrow[r,dash]&
\Phi_*^{(2)}\arrow[r,dash]&\cdots\arrow[r,dash]&
\Phi_*^{(i)}\arrow[r,"\lambda_i"]&\Phi_*^{(i+1)}
\end{tikzcd}}
satisfaisant à ceci

a) Complexe str. acyclique dans $C$, \deleted{\unsure{$\forall i\geq0$}}

b) Pour tout $i,j\geq0$, on a
\[
\operatorname{Hom}(\Phi_*^{(i)},\Phi_*^{(j)})\simeq
\left\{\begin{array}{ll}
0&\text{si }j\neq i,i+1\\
k\ (\text{engendré par id})&\text{si }j=i\\
k\ (\unsure{\text{engendré par }}\lambda_i)&\text{si }j=i+1
\end{array}\right.
\]

\deleted{\unsure{c)}} Posant
\[
\Psi_*^{(i)}=Z^{\unsure{i+1}}(\Phi_*^*)
=\operatorname{Ker}_C(\Phi^{i+1}\xrightarrow{\lambda_{\unsure{i+1}}}\Phi^{i+2})
\]
\inserted{$=\operatorname{Coker}_C(\Phi^{i-1}\xrightarrow{\lambda_{i-1}}\Phi^i)$}
\margin{\unsure{conséq. de a) b)}}
on a donc \deleted{\unsure{$\operatorname{Hom}(\Psi_i$}}
\[
\Psi^0\simeq\Phi^0
\]
\diagram{\begin{tikzcd}[column sep=large]
(*)\quad0\arrow[r,dash]&\Psi^{i-1}\arrow[r,hook,"\iota_i"]&
\Phi^i\arrow[r,two heads,"u_i"]\arrow[dr,"\lambda_i"']&
\Psi^i\arrow[r]\arrow[d,hook]&0\quad\text{ext. can.}\\
&\Phi^{i-1}\arrow[u,two heads,"u_{i-1}"]\arrow[ur,"\lambda_{i-1}"']&&\Phi^{i+1}&
\end{tikzcd}}
\[
\operatorname{Hom}(\Psi^i,\Psi^j)
\subset\operatorname{Hom}(\Phi^i,\Phi^{j+1})=
\left\{\begin{array}{ll}
0&\text{si }j\neq i\deleted{\unsure{,i+1}}\\
k&\text{si }j=i\ (\text{par id})
\end{array}\right.
\]
\[
\operatorname{Hom}(\Phi^i,\Psi^j)\deleted{\unsure{\simeq}}
\subset\operatorname{Hom}(\Phi^i,\Phi^{j+1})\simeq
\left\{\begin{array}{ll}0&\text{si }j\neq i\\
k&\text{si }j=i\ (\text{par }u_i)
\end{array}\right.
\]
\[
\operatorname{Hom}(\Psi^i,\Phi^j)
\subset\operatorname{Hom}(\Phi^i,\Phi^j)\simeq
\left\{\begin{array}{ll}
0&\text{si }j\neq i+1\ (i\neq0)\\
k&\text{si }j=i+1\ (\text{par }\iota_i)\\
k&\text{si }j=i=0\ (\text{par id})
\end{array}\right.
\]
\margin{\unsure{via composée} $\Phi^i\longrightarrow\Psi^i\longrightarrow
\Phi^{j+1}$ \illegible{3}}

d) Soit, \deleted{\unsure{$\Phi(i)$}} pour $n\in k-0$
\inserted{et $i\geq1$}, $\Phi(i,n)$
l'ext. de \deleted{\unsure{$\Phi^i$}} $\Psi^i$ par $\Psi^{i-1}$ déduite de
l'ext. $\Phi^i$ (*) par $n\operatorname{id}_{\Psi^{i-1}}$, de
sorte qu'on a un hom des suites
ex. can
\diagram{\begin{tikzcd}[column sep=large]
0\arrow[r]&\Psi^{i-1}\arrow[r,"\iota_i"]\arrow[d,"n\operatorname{id}"']&
\Phi^i\arrow[r,"u_i"]\arrow[d]&\Psi^i\arrow[r]\arrow[d,"\operatorname{id}"]&0\\
0\arrow[r]&\Psi^{i-1}\arrow[r]&\Phi(i,n)\arrow[r,dash]&\Psi^i\arrow[r,dash]&0
\end{tikzcd}}
(\uline{NB}\quad $\Phi(i,1)=\Phi^i$)
\margin{On a aussi
\diagram{\begin{tikzcd}[column sep=small,row sep=small]
0\arrow[r,dash]&\Psi^{i-1}\arrow[r,dash]\arrow[d,"\operatorname{id}"']&
\Phi(i,n)\arrow[r,dash]\arrow[d]&\Psi^i\arrow[r]\arrow[d,"n\operatorname{id}"]&0\\
0\arrow[r,dash]&\Psi^{i-1}\arrow[r]&\Phi^i\arrow[r]&\Psi^i\arrow[r]&0
\end{tikzcd}}}
Alors on a
\[
\operatorname{Hom}(\Phi(i,n),\Phi(j,m))\longrightarrow
\operatorname{Hom}(\Psi^{\unsure{i-1|i}},\Psi^j)
\]
et si $j\neq\unsure{i-1|i}$, le dernier Hom est nul,
donc si $j\neq\unsure{i-1|i}$, la structure d'extension
est respectée par tt hom $\alpha:\Phi(i,n)\longrightarrow\Phi(j,m)$
\diagram{\begin{tikzcd}[column sep=large]
0\arrow[r,dash]&\Psi^{i-1}\arrow[r,dash]\arrow[d]&\Phi(i,n)\arrow[r]\arrow[d]&
\Psi^i\arrow[r,dash]\arrow[d]&0\\
0\arrow[r,dash]&\Psi^{j-1}\arrow[r,dash]&\Phi(j,m)\arrow[r,dash]&\Psi^j\arrow[r,dash]&0
\end{tikzcd}}
Si \inserted{de plus} $j\neq i$, les flèches verticales extrêmes
sont nulles, donc on trouve
\[
\operatorname{Hom}(\Phi(i,n),\Phi(j,m))\simeq
\operatorname{Hom}(\Psi^i,\Psi^{j-1})
\]
\inserted{$\simeq k$ si $j=i+1$,\quad $0$ si $j\neq i+1$}
\margin{\unsure{via composée}
$\Phi(i,n)\longrightarrow\Psi^i\longrightarrow\Psi^{j-1}\longrightarrow\Phi(j,m)$}

\uline{Cas $j=i$}. Tout hom induit une
$p.\operatorname{id}$ sur $\Psi^{i-1}$ et $q.\operatorname{id}$ sur $\Psi^i$,
\diagram{\begin{tikzcd}[column sep=large]
0\arrow[r,dash]&\Psi^{i-1}\arrow[r,dash]\arrow[d,"p\operatorname{id}"']&
\Phi(i,n)\arrow[r,dash]\arrow[d,"\alpha"]&\Psi^i\arrow[r,dash]\arrow[d,"q\operatorname{id}"]&0\\
0\arrow[r,dash]&\Psi^{i-1}\arrow[r,dash]&\Phi(i,m)\arrow[r,dash]&\Psi^i\arrow[r,dash]&0
\end{tikzcd}}
\unsure{$p,q$ uniquement déterminés par $\alpha$.}
On a alors (pour qu'il existe $\alpha$ défini sur
$p$ et $q$) que l'élément de $\operatorname{Ext}^1(\Psi^i,\Psi^{i-1})$
défini par $q^*(e_{i,n})-p_*(e_{i,m})$ (où
$e_{i,n}$, $e_{i,m}$ désignent resp. les ext $\Phi(i,n)$ et $\Phi(i,m)$)
